\chapter*{Введение}

Современные методы компьютерной графики позволяют создавать трёхмерные сцены, отражающие объекты и процессы реального мира. Одной из практических задач является визуализация движения транспортных средств в городской среде, где важно корректно отображать геометрию, освещение и взаимодействие с элементами дороги и препятствиями.

Разработка подобных систем требует продуманной организации сцены и применения алгоритмов, обеспечивающих удаление невидимых поверхностей, расчёт освещения, текстурирование, построение теней и реализацию движения объектов. При этом необходимо поддерживать достаточную скорость работы и удобство восприятия визуализации.

\textbf{Целью курсовой работы} является разработка программного обеспечения для трёхмерного моделирования движения автомобиля в городской среде.

Для достижения поставленной цели необходимо решить следующие задачи:
\begin{enumerate}
	\item описать структуру трёхмерной сцены, включающую дороги, здания и элементы городской инфраструктуры;
	\item выбрать и реализовать алгоритмы удаления невидимых поверхностей, закрашивания и построения теней;
	\item обеспечить поддержку текстурирования и освещения от бесконечно удалённых источников света;
	\item разработать систему управления виртуальной камерой;
	\item реализовать поиск кратчайшего пути по дорожной сети и анимацию движения автомобиля от начальной до конечной точки;
	\item объединить все компоненты в единую программную систему для визуализации движения в реальном времени.
\end{enumerate}

Разрабатываемое приложение должно обеспечивать визуально достоверное отображение городской сцены с трёхмерными моделями зданий и дорог, поддержку текстур, реалистичное освещение и анимацию движения автомобиля по кратчайшему маршруту.
